% =====================================================================
% RELATÓRIO FINAL, projetos EMBRAPII
% Modelo para o relatório da última macroentrega.
% Para macroentrega intermediária, use main.tex.
%
% O relatório final é autossuficiente e serve de documento de
% transferência: quem o ler daqui a dois anos, sem acesso a mais nada,
% deve conseguir entender o projeto e continuar o trabalho. Reescreva cada
% capítulo a partir do estado final do conhecimento, usando os relatórios
% anteriores como fonte. Não cole relatórios intermediários em sequência.
%
% Compilar:    make final
% Conferir:    make verificar-final
% Entregar:    make entrega-final
% =====================================================================
\documentclass[rascunho]{macroentrega}

\setlogos{LOGO/itec_logo.png,LOGO/furg_logo.png,LOGO/embrapii_logo.png}

\dadosrelatorio{
  tipo              = final,
  projeto           = \ph{NOME DO PROJETO},
  codigo            = \ph{número no contrato e no Sisproj},
  subtitulo         = \ph{Primeira linha do subtítulo}\\\ph{segunda linha equilibrada},
  periodo           = \ph{mm/aaaa a mm/aaaa (n meses)},
  deliverable       = \ph{DN},
  parceiro          = \ph{razão social do parceiro tecnológico},
  unidade           = {Unidade EMBRAPII iTec/FURG},
  local             = \ph{Cidade},
  data              = \ph{Mês de AAAA},
  % confidencialidade = {CONFIDENCIAL},
  % codigorelatorio   = {SIGLA-RF-AAAA-NN},
  coordenador       = \ph{Nome do Coordenador},
  razaosocial       = \ph{RAZÃO SOCIAL},
  representante     = \ph{Nome do Representante},
  cargorepresentante= \ph{Cargo},
  dataassinatura    = \ph{DD de mês de AAAA},
}

\begin{document}

\capa
\clearpage\phantomsection\addcontentsline{toc}{chapter}{Equipe do Projeto}
\chapter*{Equipe do Projeto}

\orientacao{A equipe completa do projeto, organizada por função. Inclua só quem participou do projeto, e confira a grafia de cada nome.}

\begin{table}[htbp]
\centering
\estilotabela
\begin{tabularx}{\linewidth}{|P{4.2cm}|X|}
\hline
\rowcolor{reqcolor}\cab{Função} & \cab{Integrantes} \\ \hline
Coordenação & \ph{Nome Sobrenome} \\ \hline
\rowcolor{reqlight}Líderes de pesquisa & \ph{Nome Sobrenome, Nome Sobrenome} \\ \hline
Gestão & \ph{Nome Sobrenome} \\ \hline
\rowcolor{reqlight}Pesquisadores & \ph{Nome Sobrenome, Nome Sobrenome} \\ \hline
Pós-graduandos & \ph{Nome Sobrenome, Nome Sobrenome} \\ \hline
\rowcolor{reqlight}Graduandos e bolsistas & \ph{Nome Sobrenome, Nome Sobrenome} \\ \hline
\end{tabularx}
\end{table}

\begin{resumo}

\orientacao{Um parágrafo único, de 150 a 300 palavras, escrito por último, quando o documento já estabilizou. É o único trecho que todo leitor lê e o relatório inteiro comprimido, com resultado e não só com escopo. Diga, nesta ordem: o que a entrega compreende, como foi verificada, o que se obteve (com número e condição de medida) e qual o TRL declarado. Sem superlativo sobre a própria entrega.}

\ph{O que a macroentrega entregou, em uma ou duas frases.} \ph{Como a entrega foi verificada: método, ambiente e quantidade de ensaios.} \ph{Resultado principal, com número e condição de medida.} \ph{Nível de maturidade tecnológica atingido, na forma ``o projeto atinge o nível TRL~X, comprovado por Y''.}

\end{resumo}

\tableofcontents

% Em documento longo, as três listas abaixo se pagam.
\listoffigures
\listoftables
\begin{siglas}
  \item[\ph{SIGLA}] \ph{Significado, na forma por extenso.}
\end{siglas}

\chapter{Introdução}
\label{chap:introducao}

\orientacao{Diferente da introdução de uma macroentrega intermediária, a do relatório final abre no problema, porque será lida por quem não acompanhou o projeto. Cinco parágrafos, cada um com uma função:
\begin{enumerate}
    \item o problema, no domínio de aplicação, com a limitação das abordagens correntes;
    \item o desafio técnico específico e a lacuna que o projeto atacou;
    \item a hipótese ou a abordagem do projeto, em uma frase;
    \item os dados contratuais: período, duração, parceiros, arranjo de financiamento (inclusive o apoio da EMBRAPII) e o recorte de escopo, dizendo que o projeto entrega protótipo para validar métodos e algoritmos, e não produto. Essa frase custa uma linha e protege a unidade contra a expectativa de produto comercial;
    \item os pacotes de trabalho, em uma frase cada, a síntese dos resultados alcançados e a organização do documento.
\end{enumerate}}

\ph{Parágrafo 1: o problema no domínio de aplicação e a limitação das abordagens correntes.}

\ph{Parágrafo 2: o desafio técnico específico e a lacuna que o projeto atacou.}

\ph{Parágrafo 3: a hipótese ou a abordagem do projeto, em uma frase.}

O projeto \ph{NOME DO PROJETO} foi realizado pela \unidade{} com apoio da EMBRAPII, em parceria com \ph{parceiro tecnológico}, entre \ph{mm/aaaa} e \ph{mm/aaaa}. \ph{Recorte de escopo: o projeto entrega protótipo para validar métodos e algoritmos, e não produto.}

\ph{Parágrafo 5: os pacotes de trabalho, uma frase cada, e a síntese dos resultados.} O Capítulo~\ref{chap:objetivos} apresenta os objetivos; o Capítulo~\ref{chap:metodologia} descreve como o projeto foi conduzido; o Capítulo~\ref{chap:requisitos_final} consolida os requisitos; o Capítulo~\ref{chap:wp1} trata \ph{do pacote de trabalho 1}; e o Capítulo~\ref{chap:conclusao} reúne a síntese, as limitações e os trabalhos futuros.

\chapter{Objetivos}
\label{chap:objetivos}

\orientacao{Transcrição literal do contrato, sem reescrever. Se o texto não estiver à mão, deixe o marcador e peça-o à coordenação.}

\section{Objetivo Geral}

\begin{transcricao}{Anexo I, objetivo geral}
\ph{Transcrever literalmente o objetivo geral do projeto.}
\end{transcricao}

\section{Objetivos Específicos}

\begin{transcricao}{Anexo I, objetivos específicos}
\begin{enumerate}[label=\alph*)]
    \item \ph{Transcrever literalmente o primeiro objetivo específico};
    \item \ph{Transcrever literalmente o segundo objetivo específico}.
\end{enumerate}
\end{transcricao}

\chapter{Metodologia Geral}
\label{chap:metodologia}

\orientacao{Capítulo que o relatório intermediário não tem e o final precisa ter. Registra como o projeto foi conduzido: ritmo de trabalho, instâncias de reunião com o parceiro, forma de acompanhamento, composição das equipes por pacote e a lista de macroentregas com mês e conteúdo. Tem valor de auditoria, porque mostra a governança do projeto, e valor prático, porque a próxima equipe entende como o trabalho foi organizado.

Descreva o que de fato ocorreu. Se um rito foi abandonado no segundo ano, diga isso ou omita o ponto, mas não descreva um processo que não existiu.}

\section{Organização do Projeto}

\ph{Ritmo de trabalho, reuniões com o parceiro e forma de acompanhamento, como ocorreram.}

\section{Estrutura da Equipe}

\ph{Composição das equipes por pacote de trabalho e papel de cada uma.}

\section{Pacotes de Trabalho}

\ph{Os pacotes de trabalho, com o objetivo de cada um.}

\section{Macroentregas Executadas}

A Tabela~\ref{tab:macroentregas} lista as macroentregas do projeto.

\begin{table}[htbp]
\centering
\caption{Macroentregas executadas.}
\label{tab:macroentregas}
\estilotabela
\begin{tabularx}{\linewidth}{|C{2.9cm}|C{2.4cm}|X|}
\hline
\rowcolor{reqcolor}\cab{Macroentrega} & \cab{Mês} & \cab{Conteúdo, como no Anexo I} \\ \hline
\ph{M01} & \ph{mm/aaaa} & \ph{descrição} \\ \hline
\rowcolor{reqlight}\ph{M02} & \ph{mm/aaaa} & \ph{descrição} \\ \hline
\end{tabularx}
\end{table}

\chapter{Requisitos do Projeto}
\label{chap:requisitos_final}

\orientacao{Tabela consolidada com a trajetória completa do projeto, uma coluna por macroentrega, do M01 ao final. Ela serve de registro de execução para o parceiro e para a auditoria. Preencha a coluna \textbf{Verificação} nas linhas \parc{} e \done{}, com o método e o ponteiro para a evidência no relatório. Para o uso detalhado da tabela, veja o capítulo de requisitos de \texttt{main.tex}.}

Este capítulo consolida os requisitos do projeto \ph{NOME DO PROJETO} e o status de cada um ao final de cada macroentrega. O status é indicado pelas siglas \todo{} (não iniciado), \wip{} (em andamento), \parc{} (parcialmente implementado) e \done{} (concluído e verificado).

\section{Requisitos Funcionais}

\begin{tabelarequisitos}{M01,M02,M03,M04}
\req{RF-XXX-01}{\ph{O sistema deve ...}}{\wip & \wip & \done & \done}{Ensaio, Tab.~\ref{tab:wp1_resultados}}
\req{RF-XXX-02}{\ph{O sistema deve ...}}{\todo & \wip & \parc & \done}{Demonstração, Fig.~\ref{fig:wp1}}
\end{tabelarequisitos}

\section{Requisitos Não Funcionais}

\begin{tabelarequisitos}{M01,M02,M03,M04}
\req{RNF-XXX-01}{\ph{O sistema deve ...}}{\todo & \wip & \wip & \done}{Análise, Anexo~\ref{anx:a}}
\end{tabelarequisitos}

\section{Leitura da Tabela}

Dos \ph{N} requisitos do projeto, \ph{N} estão concluídos e verificados ao final da última macroentrega. \ph{Onde o avanço se concentrou, o que permanece em aberto e por qual dependência.}

\chapter{\ph{Nome do pacote de trabalho}}
\label{chap:wp1}

\orientacao{Um capítulo por pacote de trabalho (WP), ou por módulo do sistema quando o produto é um sistema único cujas partes se integram. Se as frentes tiveram equipes e métodos distintos, organize por pacote, porque o leitor procura pela frente.

Cada capítulo recompõe a frente inteira a partir do estado final do conhecimento. O resultado que foi superado aparece como etapa do percurso, não como resultado, e o experimento que fracassou aparece na discussão, não na seção de resultados.

Para o documento servir à transferência de tecnologia, algum capítulo precisa conter: lista de componentes de hardware com a função de cada um, passo a passo de montagem (quando há protótipo físico), pacotes de software com a função e a forma de comunicação, formato, volume e organização dos dados, procedimento de operação e parâmetros dos experimentos, em tabela, com unidade.}

\section{Introdução}

\ph{O que esta frente desenvolveu ao longo do projeto e como se liga às demais.}

\section{Objetivos}

\ph{Objetivos do pacote de trabalho, recortados dos objetivos específicos do projeto.}

\section{Arquitetura}

\ph{Arquitetura da solução, com a figura que a representa. Diagrama mostra intenção de projeto, não funcionamento.}

\section{Hardware}
\label{sec:wp1_hardware}

\ph{Componentes de hardware e a função de cada um; passo a passo de montagem.} A Tabela~\ref{tab:componentes} lista os componentes.

{\estilotabela
\begin{longtable}{|P{3.4cm}|P{3.8cm}|P{3.1cm}|P{3.1cm}|}
\caption{Componentes de hardware.}\label{tab:componentes}\\
\hline
\rowcolor{reqcolor}\cab{Componente} & \cab{Função} & \cab{Especificação} & \cab{Fornecedor} \\ \hline
\endfirsthead
\hline
\rowcolor{reqcolor}\cab{Componente} & \cab{Função} & \cab{Especificação} & \cab{Fornecedor} \\ \hline
\endhead
\ph{componente} & \ph{função} & \ph{valor} & \ph{fornecedor} \\ \hline
\rowcolor{reqlight}\ph{componente} & \ph{função} & \ph{valor} & \ph{fornecedor} \\ \hline
\end{longtable}}

\section{Software}

\ph{Pacotes de software, a função de cada um e como se comunicam.}

\subsection{Dados Produzidos}

\ph{Formato, volume e organização dos dados.} O dicionário de dados completo consta no Anexo~\ref{anx:b}.

\subsection{Procedimento de Operação}

\ph{Como operar o sistema, em passos que o parceiro possa seguir.}

\section{Resultados}
\label{sec:wp1_resultados}

\ph{Resultados do pacote, com condição de medida para cada número.} A Figura~\ref{fig:wp1} e a Tabela~\ref{tab:wp1_resultados} apresentam \ph{o que o leitor deve observar}.

\begin{figure}[htbp]
\centering
\figuraph[4.5cm]{figura}
\caption{\ph{Legenda que diz o que a figura mostra.}}
\label{fig:wp1}
\end{figure}

\begin{table}[htbp]
\centering
\caption{\ph{Resultados por condição de ensaio.}}
\label{tab:wp1_resultados}
\estilotabela
\begin{tabularx}{\linewidth}{|P{3.6cm}|C{2.6cm}|C{2.6cm}|X|}
\hline
\rowcolor{reqcolor}\cab{Condição} & \cab{Ensaios} & \cab{Resultado} & \cab{Observação} \\ \hline
\ph{condição} & \ph{n} & \ph{valor, unidade} & \ph{nota} \\ \hline
\end{tabularx}
\end{table}

\section{Discussão}

\ph{O que o resultado significa, sob quais condições vale e o que não demonstra.}

% Um capítulo por pacote de trabalho: copie 30_pacote_trabalho.tex.
% \include{capitulos_final/35_pacote_trabalho_2}
\chapter{Conclusão e Trabalhos Futuros}
\label{chap:conclusao}

\section{Síntese dos Resultados}

\orientacao{O que o projeto desenvolveu, organizado por pacote, em lista de resultados concretos. Cada item nomeia um artefato ou um achado, com o número que o caracteriza (quantidade de imagens, número de espécies, arquitetura usada). Feche com uma frase sobre a hipótese central: ela se confirmou, se confirmou em parte ou não se confirmou.

Declare também o nível de maturidade tecnológica atingido, na forma ``o projeto atinge o nível X, comprovado por Y'', com escala, fidelidade e ambiente, e o mesmo valor lançado no SRInfo.}

\begin{itemize}
    \item \ph{Pacote 1: artefato ou achado, com o número que o caracteriza.}
    \item \ph{Pacote 2: artefato ou achado, com o número que o caracteriza.}
\end{itemize}

\ph{A hipótese central do projeto se confirmou / se confirmou em parte / não se confirmou, porque ...} O projeto atinge o nível TRL~\ph{X}, comprovado por \ph{evidência}~\cite{abnt16290}.

\section{Limitações}

\orientacao{Uma limitação por parágrafo, com título em negrito e com causa e efeito. A primeira deve tocar o objetivo central do trabalho, quando houver limitação ali: registrá-la é o que dá autoridade ao documento, e o parceiro já a conhece.}

\paragraph{\ph{Nome da limitação}.} \ph{Causa técnica e efeito sobre o resultado.}

\paragraph{\ph{Nome da limitação}.} \ph{Causa técnica e efeito sobre o resultado.}

\section{Trabalhos Futuros}

\orientacao{Uma direção por parágrafo, com título em negrito, ligada a uma limitação ou a uma oportunidade identificada. Escreva como sugestão técnica, não como compromisso nem como proposta comercial. Se há intenção de um projeto seguinte, isso se conversa em reunião e entra em proposta, não em parágrafo de relatório.}

\paragraph{\ph{Direção}.} \ph{Sugestão técnica, ligada a uma limitação ou a uma oportunidade.}

\section{Considerações Finais}

\orientacao{Um parágrafo. Diz o que ficou demonstrado, o que o parceiro leva e o que o conjunto (conhecimento, dados, resultados experimentais) habilita. Sem superlativo e sem promessa de produto.}

\ph{Um parágrafo: o que ficou demonstrado, o que o parceiro leva e o que o conjunto habilita.}


\referencias{referencias}

\appendix
\chapter{Lista Completa de Componentes}
\label{anx:a}

\orientacao{Os anexos guardam o material que interrompe a leitura mas precisa ficar registrado. Um anexo com o inventário do que foi entregue evita a discussão sobre o que foi entregue. Cada anexo é citado no corpo do relatório. Remova os anexos que o projeto não tem.}

\ph{Lista completa de componentes, com especificação e fornecedor.}

\chapter{Estrutura e Dicionário dos Dados}
\label{anx:b}

\ph{Estrutura de diretórios dos conjuntos de dados entregues e dicionário de dados.}

\chapter{Repositórios e Versões Entregues}
\label{anx:c}

\ph{Endereço de cada repositório e a versão entregue (\emph{commit} ou \emph{tag}). Não cole trecho extenso de código: aponte o repositório. Registre a licença de bibliotecas, conjuntos de dados e modelos de terceiros quando ela impuser condição ao uso pelo parceiro.}

\chapter{Publicações e Divulgação}
\label{anx:d}

\ph{Publicações e material de divulgação produzidos no projeto, com a menção ao apoio da EMBRAPII.}

\chapter{Registros de Propriedade Intelectual}
\label{anx:e}

\ph{O que foi depositado no INPI e em que estágio está. Consulte o Núcleo de Inovação antes de descrever em detalhe qualquer resultado com potencial de proteção.}

\chapter{Material de Capacitação}
\label{anx:f}

\ph{Material de capacitação entregue ao parceiro.}


\paginaassinaturas

\end{document}
